% !PRINTED 210
\weberpage{210}
Bei der Di\"edergruppe vom Grade $n = 2m$, die wir mit $D_m$
bezeichnen wollen, haben wir nach dem vorigen Paragraphen
zwei Systeme von je $m$ conjugirten zweizähligen Polen und ein
System von zwei conjugirten $m$-zähligen Polen. Die letzteren
müssen also nach dem Satze \S.\,53 die Pole einer Substitution
sein. Wenn wir sie mit $0$ und $\infty$ zusammenfallen lassen, so erhalten
 wir die Grundform zweiten Grades:
\[
f_1 = x_1 x_2.
\]
Dies kann aber nur dann eine invariante Form der Gruppe
sein, wenn alle Substitutionen in einer der beiden Formen
\[
\begin{pmatrix} \alpha, & 0 \\ 0, & \delta \end{pmatrix}, \quad \begin{pmatrix} 0, & \beta \\ \gamma, & 0 \end{pmatrix}
\]
enthalten sind, und hierin müssen $\alpha:\delta$ und $-\beta:\gamma$ $m^{\text{te}}$ Einheitswurzeln
 sein. Wir bekommen also die Substitutionen der Gruppe,
wenn $\varepsilon$ eine primitive $m^{\text{te}}$ Einheitswurzel ist:
\[
x, \varepsilon x, \varepsilon^2 x, \ldots, \varepsilon^{m-1} x
\]
\[
\frac{1}{x}, \frac{\varepsilon}{x}, \frac{\varepsilon^2}{x}, \cdots, \frac{\varepsilon^{m-1}}{x},
\]
d.\,h.
\weq{5}{D_m = \begin{pmatrix} \varepsilon^r, & 0 \\ 0, & 1 \end{pmatrix}, \begin{pmatrix} 0, & \varepsilon^r \\ 1, & 0 \end{pmatrix},}
oder, mit der Determinante $+1$ dargestellt:
\weq{6}{D_m = \begin{pmatrix} \varepsilon^{1/2\,r}, & 0 \\ 0, & \varepsilon^{-1/2\,r} \end{pmatrix}, \begin{pmatrix} 0, & i\varepsilon^{1/2\,r} \\ i\varepsilon^{-1/2\,r}, & 0 \end{pmatrix}, \quad r = 0, 1, \ldots, m-1.}
Man sieht sofort, dass diese Substitutionen wirklich eine
Gruppe bilden.

Ausser $0$ und $\infty$ haben wir hier noch die aus den Gleichungen
\[
x = \frac{\varepsilon^h}{x}
\]
hervorgehenden Pole $\pm \varepsilon^{1/2\,h}$, aus denen man noch die beiden
Grundformen $m^{\text{ten}}$ Grades
\weq{7}{f_2 = x_1^m + x_2^m, \quad f_3 = x_1^m - x_2^m}
erhält.

Hieraus ersieht man, dass die Di\"edergruppen und die cyklischen
 Gruppen ganz von einander verschieden sind, da ihre
Grundformen verschiedene Grade haben.
